\sectioncentered*{Заключение}
\addcontentsline{toc}{section}{Заключение}

В рамках работы разработана информационно-поисковая система, реализующая интерфейс с пользователем на основе вероятностной модели поиска Okapi BM25. Система разворачивается в локальной вычислительной сети средствами Docker Compose и предназначена для работы с документами на русском языке, что соответствует варианту~5 задания.

В ходе выполнения работы решены следующие задачи:

\begin{itemize}
	\item спроектирована и реализована трёхуровневая архитектура: веб-интерфейс на React, серверная часть на Python 3.11 (FastAPI) и уровень данных на PostgreSQL 16;
	\item спроектирована база данных из восьми таблиц: документы, словарь, инвертированный индекс, журналы поиска, эталонные запросы, разметка релевантности, прогоны оценки и учёт миграций;
	\item реализованы алгоритмы предобработки русского текста (токенизация, лемматизация, фильтрация стоп-слов), построения инвертированного индекса, вычисления весов Робертсона -- Спарк Джонс и ранжирования по BM25;
	\item реализован интеллектуальный интерфейс: автодополнение, исправление опечаток, расширение запроса синонимами, сниппеты с подсветкой и список слов запроса, найденных в документе, встроенная справка;
	\item реализована загрузка документов форматов TXT, MD, PDF, DOCX, HTML, RTF, CSV и LOG;
	\item реализована программная оценка качества: точность, полнота, F-меры, P@k, MAP, nDCG и 11-точечная кривая полноты-точности, с выводом таблиц и графиков;
	\item обеспечено воспроизводимое развёртывание Docker Compose с автоматическим применением миграций и загрузкой корпуса.
\end{itemize}

Тестирование подтвердило работоспособность предобработки, индексирования, поиска, интеллектуального интерфейса, загрузки файлов и оценки качества. На коллекции из 18 разнородных русскоязычных документов и 18 эталонных запросов модель BM25 показала среднюю точность 0,736, полноту 0,981 и F1-меру 0,798, превзойдя базовую векторную модель TF-IDF.

Перспективными направлениями улучшения являются расширение корпуса, учёт словосочетаний, адаптивный порог отсечения выдачи, замена ручного тезауруса на векторные представления и использование подлинной обратной связи пользователей.

Поставленная цель достигнута: система работоспособна, соответствует вероятностной стратегии поиска и может использоваться как основа для дальнейших исследований в области информационного поиска.

\begin{thebibliography}{9}
	\addcontentsline{toc}{section}{Cписок использованных источников}

	\bibitem{manning}
	Маннинг К., Рагхаван П., Шютце Х. Введение в информационный поиск. "--- М. : Вильямс, 2011. "--- 528~с.

	\bibitem{robertson}
	Robertson S., Zaragoza H. The Probabilistic Relevance Framework: BM25 and Beyond // Foundations and Trends in Information Retrieval. "--- 2009. "--- Vol.~3, No.~4. "--- P.~333--389.

	\bibitem{rsj}
	Robertson S. E., Sparck Jones K. Relevance weighting of search terms // Journal of the American Society for Information Science. "--- 1976. "--- Vol.~27, No.~3. "--- P.~129--146.

	\bibitem{romip}
	Российский семинар по оценке методов информационного поиска (ROMIP). Метрики качества поиска. "--- 2004.

	\bibitem{pymorphy}
	Korobov M. Morphological Analyzer and Generator for Russian and Ukrainian Languages // Analysis of Images, Social Networks and Texts. "--- 2015. "--- P.~320--332.

	\bibitem{fastapi}
	FastAPI documentation [Электронный ресурс]. "--- Режим доступа: \url{https://fastapi.tiangolo.com/}.

	\bibitem{postgres}
	PostgreSQL 16 Documentation [Электронный ресурс]. "--- Режим доступа: \url{https://www.postgresql.org/docs/16/}.

	\bibitem{react}
	React documentation [Электронный ресурс]. "--- Режим доступа: \url{https://react.dev/}.
\end{thebibliography}
